\section{Introduction}
\label{sec:introduction}

\IEEEPARstart{D}{igital} forensics and incident response face an acute evidentiary bottleneck. Forensic laboratories routinely seize enterprise storage arrays and personal devices containing dozens of terabytes of unstructured text, volatile memory images, and heterogeneous audit logs. Classical investigative methodologies formalized by Carrier~\cite{carrier2005} and Casey~\cite{casey2011} depend on deterministic string indexing, regular expressions, and byte-level file carving. Although these foundational techniques preserve strict mathematical integrity, they scale poorly against multi-stage intrusions where decisive contextual indicators span millions of dispersed operating system events.

Transformer-based Large Language Models (LLMs) and Generative Artificial Intelligence (GenAI) offer a substantial shift in investigative triage capacity. Practitioners and researchers apply proprietary frontier models alongside fine-tuned local architectures to summarize multi-source logs, classify suspicious system call sequences, triage volatile memory dumps, and draft structured case notes~\cite{scanlon2023_digital, scanlon2023_chatgpt, oh2024_volgpt}. Controlled evaluations report marked gains in initial triage speed, multilingual chat inspection, and documentation workflows~\cite{michelet2024_chatgpt, kim2025_digital, wickramasekara2025_exploring}. Concurrently, the rapid proliferation of local and agentic AI assistants on consumer endpoints has created an inverse investigative challenge: recovering conversational histories, tool invocations, and Model Context Protocol (MCP) traces from suspect devices~\cite{walker2024_forensic, xu2026_recovering, satter2026_with}.

Computational acceleration introduces direct legal and technical friction. Digital forensic science operates under rigorous procedural standards governing acquisition, preservation, analysis, and method validation, notably ISO/IEC 27037~\cite{iso27037_2012}, ISO/IEC 27041~\cite{iso27041_2015}, ISO/IEC 27042~\cite{iso27042_2015}, and NIST Special Publication 800-86~\cite{nist_sp800_86}. In judicial proceedings, expert methodology faces strict scrutiny under the Daubert standard~\cite{daubert1993}, Federal Rule of Evidence (FRE) 702~\cite{fre702_2023}, continental criminal procedure codes~\cite{portugal_cpp}, and the high-risk AI obligations established by the European Union Artificial Intelligence Act~\cite{eu_ai_act_2024}. Commercial LLMs function as opaque, non-deterministic services with undisclosed training corpora. Transmitting unredacted evidentiary artifacts to external cloud APIs compromises statutory data sovereignty and chain-of-custody accounting~\cite{sharma2025_forensicllm, eu_led_2016}. Recent empirical evidence also reveals a subtle epistemic hazard: rather than merely mirroring leading investigator hypotheses, frontier models exhibit prompt-dependent error rates and fail to challenge unsupported premises embedded inside examiner queries~\cite{trivedi2026_when}. This compromises evidentiary reliability.

Beyond reliability, LLM integration expands the adversarial attack surface of forensic laboratories. Suspect files, event logs, and command-line records ingested by automated triage pipelines can embed indirect prompt injection payloads. An adversary anticipating post-incident examination can plant natural language directives inside passive system artifacts, attempting to steer the generative triage engine into suppressing malicious process identifiers or reclassifying hostile commands as benign administration~\cite{chernyshev2026_large, kim2026_evidenceborne}. Forensic engineering must reconcile semantic automation with deterministic verification and strict defensive isolation.

Despite rapid academic output, the literature remains scattered across digital forensics journals, software security venues, and legal jurisprudence periodicals. Existing secondary studies offer valuable early perspectives~\cite{scanlon2023_chatgpt, wickramasekara2025_exploring, chernyshev2026_large, jelodar2026_large, scanlon2026_investigation}, yet none simultaneously combine a reproducible PRISMA 2020 protocol, Crossref--OpenAlex bibliometric triangulation, decoupled orthogonal taxonomy across both forensic directions, hardware determinism modeling, European high-risk AI / LED statutory analysis, and a public replication package. Table~\ref{tab:survey_comparison} positions this systematic review against existing secondary surveys in the domain.

\begin{table*}[!t]
\caption{Positioning of This Systematic Review Against Prior Secondary Surveys and Roadmaps in LLM Digital Forensics}
\label{tab:survey_comparison}
\centering
\scriptsize
\setlength{\tabcolsep}{3.5pt}
\begin{tabular}{p{3.4cm}cccccccc}
\toprule
\textbf{Secondary Study (Year)} & \textbf{PRISMA} & \textbf{Biblio. +} & \textbf{Dual Axis} & \textbf{QA1--QA4} & \textbf{Determinism} & \textbf{EU AI Act /} & \textbf{Anti-} & \textbf{Public} \\
 & \textbf{2020} & \textbf{OpenAlex} & \textbf{(For/Of)} & \textbf{Audit} & \textbf{Manifest} & \textbf{LED / CPP} & \textbf{Forensics} & \textbf{Repo} \\
\midrule
Scanlon et al.~\cite{scanlon2023_chatgpt} (2023) & -- & -- & -- & -- & -- & -- & Part. & -- \\
Wickramasekara et al.~\cite{wickramasekara2025_exploring} (2025) & -- & -- & -- & -- & -- & Part. & Part. & -- \\
Chernyshev et al.~\cite{chernyshev2026_large} (2026) & -- & -- & Part. & -- & -- & Part. & \checkmark & -- \\
Jelodar et al.~\cite{jelodar2026_large} (2026) & -- & -- & -- & -- & -- & -- & Part. & -- \\
Scanlon~\cite{scanlon2026_investigation} (2026) & -- & -- & \checkmark & -- & -- & Part. & -- & -- \\
\midrule
\textbf{This Review (2026)} & \textbf{\checkmark} & \textbf{\checkmark} & \textbf{\checkmark} & \textbf{\checkmark} & \textbf{\checkmark} & \textbf{\checkmark} & \textbf{\checkmark} & \textbf{\checkmark} \\
\bottomrule
\multicolumn{9}{p{16.5cm}}{\tiny Legend: \checkmark\ = full systematic/formal treatment; Part. = partial or high-level discussion; -- = not addressed. Dual Axis separates LLMs applied \textit{for} DFIR from forensics \textit{of} LLM/agentic systems.}
\end{tabular}
\end{table*}

To close this gap, this study conducts a systematic review and bibliometric analysis adhering to the PRISMA 2020 protocol~\cite{page2021prisma}. We retrieve, screen, and synthesize peer-reviewed publications indexed via the Crossref REST API from January 2021 through November 10, 2026 ($N=44$), explicitly distinguishing LLMs applied \textit{to} forensic pipelines from the forensic investigation \textit{of} LLM systems. Four research questions guide this investigation:
\begin{itemize}
    \item \textbf{RQ1 (Bibliometric and Structural Profile):} How has scholarly output evolved across time, venues, geographic clusters, and orthogonal structural axes (forensic direction, deployment locality, and agentic orchestration)?
    \item \textbf{RQ2 (Operational Capabilities and Forensic Targets):} Across which concrete operational domains (triage, incident response, malware analysis, memory forensics, and forensics of LLM systems) are LLMs empirically validated?
    \item \textbf{RQ3 (Custody, Verifiability, and Legal Compliance):} What cryptographic manifests, reproducibility controls, and local execution architectures satisfy ISO/IEC 27037/27042, Daubert, FRE 702(d), EU data protection regimes (GDPR and Directive 2016/680), and the EU AI Act?
    \item \textbf{RQ4 (Anti-Forensic Threats and Defensive Hardening):} Which adversarial mechanisms, including evidence-borne prompt injection, behavioral malware evasion, context flooding, and chronological distortion, compromise LLM-assisted pipelines, and what architectural defenses counter them?
\end{itemize}

Section~\ref{sec:methodology} documents the PRISMA 2020 search protocol, exact query suite, eligibility boundaries, and quality appraisal framework. Section~\ref{sec:bibliometrics} presents quantitative bibliometric findings and orthogonal taxonomic metrics answering RQ1. Section~\ref{sec:use_cases} synthesizes operational implementations and LLM-artifact investigations addressing RQ2. Section~\ref{sec:custody_crypto} evaluates cryptographic integrity, deterministic replay, and multi-jurisdictional admissibility under RQ3. Section~\ref{sec:anti_forensics} examines anti-forensic threat vectors and hardening strategies under RQ4. Section~\ref{sec:discussion} analyzes structural literature gaps, future research directions, and threats to validity, and Section~\ref{sec:conclusion} concludes the paper.
